% Seção: Trabalhos relacionados  (orçamento: 1 a 1,5 página)
%
% Formato combinado com o autor: quatro trabalhos com um parágrafo cada (SoccerRAG, SPORTSQL, Tzikas, LOFT) e os
% demais em uma frase, mais a tabela de posicionamento. Esta seção só posiciona: comparação com os resultados
% (por exemplo, o acerto de Tzikas contra o da Busca vetorial) fica na Discussão.
% Fora daqui: NLGraph, GraphRAG (Edge), StructGPT, Pan et al. e Pappalardo; Duch, Buldú e Sui ficam para a Discussão,
% se um resultado pedir. "Pergunta encadeada" é definida na Seção 5.3; aqui a definição curta vem com a remissão.
% Números de outros trabalhos vão por extenso (o lint só aceita números do benchmark como macro).

\section{Trabalhos relacionados}\label{sec:trabalhos}

Os primeiros conjuntos de avaliação de LLMs no esporte medem o que o modelo sabe, sem dados de uma partida. O
SportQA~\cite{xia2024sportqa} reúne mais de setenta mil perguntas de múltipla escolha, de fatos históricos a
cenários de jogo; os modelos avaliados foram bem nos fatos básicos e ficaram atrás de especialistas humanos nos
cenários. O SoccerAgent~\cite{rao2025socceragent} divide as perguntas entre vários agentes que consultam uma base
de conhecimento sobre jogadores, times, árbitros e estádios, também com perguntas de múltipla escolha.

Os trabalhos mais próximos ligam o modelo a dados de partidas. O SoccerRAG~\cite{strand2024soccerrag} guarda em
um banco SQLite as anotações de eventos do SoccerNet, como gols, cartões e faltas, os comentários transcritos das
transmissões e os dados dos jogadores. O modelo extrai e valida os nomes citados na pergunta, e um agente escreve
a consulta SQL. A avaliação usou vinte perguntas de exemplo e comparou seis configurações do próprio sistema em
dez delas, todas com o modelo escrevendo SQL.

O SPORTSQL~\cite{martinez2025sportsql} traduz perguntas sobre a liga inglesa em consultas SQL sobre um banco
montado com os dados do Fantasy Premier League e indexado no tempo. Acompanha um conjunto de mais de mil e
setecentas perguntas, cada uma com a consulta SQL e a resposta de referência, o que mostra que a conferência por
resposta exata é viável no esporte. Os dados, porém, são estatísticas de jogadores e times, e não o registro de
cada ação da partida.

Tzikas et al.~\cite{tzikas2025dynamic} usam eventos de cinco rodadas da liga alemã, com posição no campo, tempo e
jogador. Cada evento vira uma frase montada a partir de um molde de texto, e a resposta vem de recuperação
aumentada por geração, com reordenação dos trechos recuperados. As trinta perguntas são de três tipos (quando,
onde e quem), e a conferência é feita por um LLM, que julga se a resposta tem o mesmo sentido do gabarito. Os
autores relatam erros desse julgamento e apontam a recuperação iterativa como caminho para as perguntas que
exigem juntar informação espalhada pela partida.

Fora do esporte, a comparação entre formas de recuperação já foi estudada em dados gerais.
Li et al.~\cite{li2024ragorlongcontext} compararam recuperação e contexto longo em conjuntos públicos: com
recursos suficientes, o contexto longo teve desempenho médio maior, e a recuperação, custo bem menor. Liu
et al.~\cite{liu2024lostinthemiddle} mostraram que o desempenho cai quando a informação relevante está no meio de
um contexto longo. Han et al.~\cite{han2025ragvsgraphrag} compararam a recuperação de trechos de texto com a
recuperação ao longo de um grafo, sob um mesmo protocolo, e cada forma se saiu melhor em tarefas diferentes.
Quando os dados estão em um banco de grafos, o modelo também pode escrever a consulta, mas LLMs ainda erram ao
gerar consultas Cypher~\cite{ozsoy2025text2cypher}.

O LOFT~\cite{lee2024loft} é o que mais se aproxima da comparação feita aqui. O conjunto testa se um modelo que
recebe o corpus inteiro no contexto substitui sistemas de recuperação, de recuperação aumentada por geração e de
SQL. O contexto longo teve desempenho comparável ao dos sistemas de recuperação, mas teve dificuldade com o
raciocínio composicional das tarefas do tipo SQL, isto é, as que combinam várias partes dos dados.

\vspace{2mm}
\noindent\begin{minipage}{\columnwidth}
\centering\footnotesize\renewcommand{\arraystretch}{1.15}\setlength{\tabcolsep}{3pt}
\begin{tabular}{>{\raggedright\arraybackslash}p{0.25\columnwidth}>{\raggedright\arraybackslash}p{0.235\columnwidth}>{\raggedright\arraybackslash}p{0.21\columnwidth}>{\raggedright\arraybackslash}p{0.21\columnwidth}} % lint: ignore
\toprule
Trabalho & Dados & Formas de recuperação & Conferência \\
\midrule
SportQA \cite{xia2024sportqa} & nenhum (memória do modelo) & nenhuma & múltipla escolha \\
SoccerAgent \cite{rao2025socceragent} & base de conhecimento & uma (agentes) & múltipla escolha \\
SoccerRAG \cite{strand2024soccerrag} & anotações e comentários de partidas & uma (SQL), em seis configurações & feita pelos autores \\
SPORTSQL \cite{martinez2025sportsql} & estatísticas de jogadores e times & uma (SQL) & resposta de referência \\
Tzikas et al. \cite{tzikas2025dynamic} & eventos de partidas & uma (recuperação vetorial) & LLM como juiz \\
Li et al. \cite{li2024ragorlongcontext} & textos gerais & duas (recuperação e contexto longo) & métricas dos conjuntos \\
LOFT \cite{lee2024loft} & textos e tabelas gerais & contexto longo contra sistemas especializados & métricas dos conjuntos \\
\bottomrule
\end{tabular}
\captionof{table}{Trabalhos relacionados: dados, formas de recuperação e conferência da resposta.}\label{tab:posicionamento}
\end{minipage}\par\smallskip

A Tabela~\ref{tab:posicionamento} resume esses trabalhos. No esporte, cada trabalho avalia um sistema com uma
única forma de recuperação; fora dele, as formas de recuperação são comparadas, mas sem dados esportivos. Não foram
encontrados trabalhos que comparem formas de recuperação em perguntas encadeadas sobre eventos de futebol, isto é,
perguntas cuja resposta exige ligar eventos entre si (Seção~\ref{sec:perguntas}), com conferência por código
contra um gabarito calculado dos dados brutos.
